% GENERATED FILE. Edit canonical lessons, then run npm run generate:books.
% canonical-source-hash: fnv1a64:db0388842aadbaf5
% canonical-lessons: ML-C161-mannu, ML-C161-tennu, ML-C161-vala2, ML-C161-mulla, ML-C161-tamara

\chapter{Soil, Coconut palm, Banana plant, Jasmine, Lotus}
\label{ch:ml-soil-coconut-palm-banana-plant-jasmine-lotus}
\hlchaptermodality{161}

\begin{chapteropening}
\textbf{By the end of this chapter:} \emph{I can say soil, a coconut palm, a banana plant, jasmine and a lotus.}

Complete the last lesson of chapter 161: i can say soil, a coconut palm, a banana plant, jasmine and a lotus.
\end{chapteropening}

\section[\texorpdfstring{\ml{മണ്ണ്} (maṇṇŭ)}{മണ്ണ് (maṇṇŭ)}]{\ml{മണ്ണ്} (maṇṇŭ) — soil, earth}

\label{lesson:ML-C161-mannu}



\begin{quote}
Before the new one: say the Malayalam for a waterfall, then the Malayalam for a stone.

Say \textquotedblleft{}midnight.\textquotedblright{} (\emph{Pātirā} — \emph{pāti}, \textquotedblleft{}half,\textquotedblright{} plus night.)
\end{quote}

\subsection*{You'll want to know: \ml{മണ്ണ്}}
\textbf{\ml{മണ്ണ്}} — \emph{maṇṇŭ} — \textquotedblleft{}soil, earth\textquotedblright{}.

\begin{grammarlens}[title={What you've built}]
One more word for this chapter.
\end{grammarlens}

\subsection*{Your turn}
\begin{itemize}
\raggedright
  \item \emph{Say it:} \emph{maṇṇŭ}
  \item \emph{Say it:} \emph{maṇṇŭ}, once more
  \item \emph{Say it:} say \emph{maṇṇŭ}
  \item \emph{Recall:} say the Malayalam for a waterfall, then the Malayalam for a stone, then say \emph{maṇṇŭ} again
\end{itemize}

\begin{tcolorbox}[breakable,colback=teal!4,colframe=teal!35!black,title={Before you move on}]
What does \ml{മണ്ണ്} mean? (\textquotedblleft{}Soil, earth\textquotedblright{}.) Say it once more.
\end{tcolorbox}
\section[\texorpdfstring{\ml{തെങ്ങ്} (teṅṅŭ)}{തെങ്ങ് (teṅṅŭ)}]{\ml{തെങ്ങ്} (teṅṅŭ) — a coconut palm}

\label{lesson:ML-C161-tennu}



\begin{quote}
Before the new one: say the Malayalam for a stone, then the Malayalam for soil.
\end{quote}

\subsection*{You'll want to know: \ml{തെങ്ങ്}}
\textbf{\ml{തെങ്ങ്}} — \emph{teṅṅŭ} — \textquotedblleft{}a coconut palm\textquotedblright{}.

\begin{grammarlens}[title={What you've built}]
One more word for this chapter.
\end{grammarlens}

\subsection*{Your turn}
\begin{itemize}
\raggedright
  \item \emph{Say it:} \emph{teṅṅŭ}
  \item \emph{Say it:} \emph{teṅṅŭ}, once more
  \item \emph{Say it:} say \emph{teṅṅŭ}
  \item \emph{Recall:} say the Malayalam for a stone, then the Malayalam for soil, then say \emph{teṅṅŭ} again
\end{itemize}

\begin{tcolorbox}[breakable,colback=teal!4,colframe=teal!35!black,title={Before you move on}]
What does \ml{തെങ്ങ്} mean? (\textquotedblleft{}A coconut palm\textquotedblright{}.) Say it once more.
\end{tcolorbox}
\section[\texorpdfstring{\ml{വാഴ} (vāḻa)}{വാഴ (vāḻa)}]{\ml{വാഴ} (vāḻa) — a banana plant}

\label{lesson:ML-C161-vala2}



\begin{quote}
Before the new one: say the Malayalam for soil, then the Malayalam for a coconut palm.

Say \textquotedblleft{}noon.\textquotedblright{} (\emph{Ucca}.) Is it a \emph{madhya}, \textquotedblleft{}middle,\textquotedblright{} compound like the Telugu word? (No: a different root.)
\end{quote}

\subsection*{You'll want to know: \ml{വാഴ}}
\textbf{\ml{വാഴ}} — \emph{vāḻa} — \textquotedblleft{}a banana plant\textquotedblright{}.

\begin{grammarlens}[title={What you've built}]
One more word for this chapter.
\end{grammarlens}

\subsection*{Your turn}
\begin{itemize}
\raggedright
  \item \emph{Say it:} \emph{vāḻa}
  \item \emph{Say it:} \emph{vāḻa}, once more
  \item \emph{Say it:} say \emph{vāḻa}
  \item \emph{Recall:} say the Malayalam for soil, then the Malayalam for a coconut palm, then say \emph{vāḻa} again
\end{itemize}

\begin{tcolorbox}[breakable,colback=teal!4,colframe=teal!35!black,title={Before you move on}]
What does \ml{വാഴ} mean? (\textquotedblleft{}A banana plant\textquotedblright{}.) Say it once more.
\end{tcolorbox}
\section[\texorpdfstring{\ml{മുല്ല} (mulla)}{മുല്ല (mulla)}]{\ml{മുല്ല} (mulla) — jasmine}

\label{lesson:ML-C161-mulla}



\begin{quote}
Before the new one: say the Malayalam for a coconut palm, then the Malayalam for a banana plant.
\end{quote}

\subsection*{You'll want to know: \ml{മുല്ല}}
\textbf{\ml{മുല്ല}} — \emph{mulla} — \textquotedblleft{}jasmine\textquotedblright{}.

\begin{grammarlens}[title={What you've built}]
One more word for this chapter.
\end{grammarlens}

\subsection*{Your turn}
\begin{itemize}
\raggedright
  \item \emph{Say it:} \emph{mulla}
  \item \emph{Say it:} \emph{mulla}, once more
  \item \emph{Say it:} say \emph{mulla}
  \item \emph{Recall:} say the Malayalam for a coconut palm, then the Malayalam for a banana plant, then say \emph{mulla} again
\end{itemize}

\begin{tcolorbox}[breakable,colback=teal!4,colframe=teal!35!black,title={Before you move on}]
What does \ml{മുല്ല} mean? (\textquotedblleft{}Jasmine\textquotedblright{}.) Say it once more.
\end{tcolorbox}
\section[\texorpdfstring{\ml{താമര} (tāmara)}{താമര (tāmara)}]{\ml{താമര} (tāmara) — a lotus}

\label{lesson:ML-C161-tamara}



\begin{quote}
Before the new one: say the Malayalam for a banana plant, then the Malayalam for jasmine.

\emph{Maṇi}: a bell, or an hour? (Both; Tamil dictionaries treat them as two look-alike words.)
\end{quote}

\subsection*{You'll want to know: \ml{താമര}}
\textbf{\ml{താമര}} — \emph{tāmara} — \textquotedblleft{}a lotus\textquotedblright{}.

\begin{grammarlens}[title={What you've built}]
One more word for this chapter.
\end{grammarlens}

\subsection*{Your turn}
\begin{itemize}
\raggedright
  \item \emph{Say it:} \emph{tāmara}
  \item \emph{Say it:} \emph{tāmara}, once more
  \item \emph{Say it:} say \emph{tāmara}
  \item \emph{Recall:} say the Malayalam for a banana plant, then the Malayalam for jasmine, then say \emph{tāmara} again
\end{itemize}

\begin{tcolorbox}[breakable,colback=teal!4,colframe=teal!35!black,title={Before you move on}]
What does \ml{താമര} mean? (\textquotedblleft{}A lotus\textquotedblright{}.) Say it once more.
\end{tcolorbox}
